% Directly authored companion to gamma_optimizer_access_note.md.
\documentclass[11pt]{article}
\usepackage[T1]{fontenc}
\usepackage{amsmath,amssymb,amsthm,newtxtext,newtxmath}
\usepackage[letterpaper,margin=0.78in]{geometry}
\usepackage{graphicx,microtype,xcolor,caption,fancyhdr,tabularx,pdflscape,placeins}
\usepackage[colorlinks=true,linkcolor=blue!40!black,urlcolor=blue!50!black]{hyperref}
\definecolor{ink}{HTML}{263442}
\DeclareMathOperator{\sech}{sech}
\DeclareMathOperator{\sign}{sign}
\DeclareMathOperator{\diag}{diag}
\theoremstyle{definition}\newtheorem{theorem}{Theorem}
\graphicspath{{results/checkpoint_D_optimizers/expD36_frozen_gamma_probe/full_sweep/refinements/gamma_optimizer_access/}}
\hypersetup{pdftitle={How gamma controls access to a target during readout training}}
\setlength{\parindent}{0pt}\setlength{\parskip}{5pt plus 1pt}
\setlength{\emergencystretch}{2em}\setlength{\headheight}{14pt}
\captionsetup{font=small,labelfont=bf,justification=raggedright,singlelinecheck=false}
\pagestyle{fancy}\fancyhf{}
\fancyhead[L]{\small\color{ink}Gamma and access to a target}
\fancyhead[R]{\small Mechanism, GD, and Adam}
\fancyfoot[C]{\small\thepage}\renewcommand{\headrulewidth}{0.3pt}
\begin{document}
\thispagestyle{plain}
{\color{ink}\fontsize{23}{26}\selectfont\bfseries
How gamma controls access to a target\\during readout training\par}
\vspace{8pt}
Small slopes can make an attainable target expensive to learn. In our
frozen-feature tanh model, the same target reaches 1\% relative error after
\textbf{15.8 million GD updates at gamma 8}, compared with \textbf{16,013
updates at gamma 64}. Adam also struggles at gamma 8: error remains above
1\% after 200,000 updates. Centers, target, readout coordinates, and zero
initialization are held fixed.

The explanation has three parts. Gamma determines how strongly each
feature smooths rapid variation. This changes the kernel's response to
patterns in the target, and the measured target-weighted spectrum reveals
their relative learning rates. Classical GD dynamics then predict the
observed delay; Adam tests empirically how the difficulty persists under
adaptive updates. We follow one representative target, then check other
targets and optimizer settings.

\textbf{Notation.} Eigenvalues are denoted by $\mu$ throughout.
\begin{center}\small
\begin{tabularx}{\linewidth}{@{}p{0.34\linewidth}X@{}}
\hline Symbol & Meaning\\\hline
$\gamma$ & Common slope of frozen tanh features\\
$J_\gamma$, $K_\gamma=J_\gamma J_\gamma^T$ & Sample-normalized feature matrix and readout kernel\\
$y$, $r_n$, $e(n)=\|r_n\|/\|y\|$ & Target samples, residual, and relative error\\
$\mu_i$, $u_i$, $\rho_i=\mu_i/\mu_1$ & Eigenvalues, unit eigenvectors, relative eigenvalues\\
$p_i=|u_i^Ty|^2/\|y\|^2$ & Fraction of target energy in an eigenvector\\
$q_\gamma(v)=v^TK_\gamma v$ & Kernel response to fixed unit sample pattern $v$\\\hline
\end{tabularx}
\end{center}

\section{The target is attainable; the question is how quickly}
The readout model is
\[
\widehat f(x)=w_0+\sum_{j=1}^{W}w_j\tanh\!\left(\gamma(x-c_j)\right).
\]
Only coefficients $w_j$ are trained. Centers $c_j$ and common slope
$\gamma$ remain fixed during each run. We compare
$\gamma\in\{8,12,16,64\}$ on the same uniform centers and samples.
The primary target is
\[
f(x)=\sin(2\pi x)+\tfrac12\sin(6\pi x)+\tfrac14\sin(10\pi x),
\qquad x\in[-1,1].
\]
There are 559 tanh features plus bias, spacing $h=1/256$, and
$m=8{,}193$ training samples. Normalize features and target by $\sqrt m$:
\[
J_\gamma[a,j]=\frac{\tanh(\gamma(x_a-c_j))}{\sqrt m},\qquad
J_\gamma[a,0]=\frac1{\sqrt m},\qquad y_a=\frac{f(x_a)}{\sqrt m}.
\]
Then $\tfrac12\|J_\gamma w-y\|^2$ is half mean-squared error.
All four GD runs actually reach 1\% error, establishing attainability at
the studied accuracy, without requiring exact representation everywhere
on the interval.

The kernel measures how residual patterns couple to available features.
For a unit vector $v$ on the sample grid,
\[
q_\gamma(v)=v^TK_\gamma v=\|J_\gamma^Tv\|^2.
\]
This is total squared correlation of the dictionary with that pattern.
For an eigenvector it equals the eigenvalue. For a general direction,
\[
q_\gamma(v)=\sum_i\mu_i(\gamma)|u_i(\gamma)^Tv|^2.
\]
The same quadratic expression defines $q_\gamma(v)$ for arbitrary $v$;
unit normalization makes directional responses comparable.
We can track a physically fixed pattern across gamma while allowing
the eigenvectors to change.

\section{Gamma explicitly controls the kernel response}
\subsection*{Smaller gamma smooths more strongly}
A tanh transition is a smoothed step:
\[
\tanh(\gamma\,\cdot)=s_\gamma*\sign,\qquad
s_\gamma(x)=\frac\gamma2\sech^2(\gamma x).
\]
Convolution averages translated steps against the density $s_\gamma$.
Small gamma spreads this average more widely. Its Fourier multiplier is
\[
M_\gamma(\omega)=\frac z{\sinh z},\qquad
z=\frac{\pi|\omega|}{2\gamma},\qquad M_\gamma(0)=1.
\]
It is close to one for slow variation and decays rapidly when
$|\omega|/\gamma$ is large. Increasing gamma admits more rapid variation.
The squared multiplier enters the kernel because it pairs two smoothed
features. This describes original tanh features; the step explains their
construction.

\subsection*{An explicit statement for the finite tanh dictionary}
Inputs and centers stay fixed. For $d_{ab}=x_a-x_b$, define the analytic
center-integral contribution and finite correction:
\[
B_\gamma[a,b]=\frac{W+1}{m}-\frac2{hm}d_{ab}\coth(\gamma d_{ab}),
\qquad R_\gamma=K_\gamma-B_\gamma.
\]
Use the continuous value $d\coth(\gamma d)=1/\gamma$ at zero.
$R_\gamma$ accounts for the finite center interval and replacement of
its integral by a discrete sum. These effects belong to the actual model.
The proof gives their explicit integral expression; the example measures
their size.

\begin{theorem}[Gamma-dependent response of a finite tanh kernel]
For consecutive uniformly spaced centers, fixed input samples, slopes
$\Gamma>\gamma>0$, and every real sample vector $v$, define
\[
V_v(\omega)=\sum_{a=1}^{m}v_a e^{-i\omega x_a}.
\]
Then
\begin{align*}
q_\Gamma(v)-q_\gamma(v)
&=\mathcal A_{\Gamma,\gamma}(v)+v^T(R_\Gamma-R_\gamma)v,\\
\mathcal A_{\Gamma,\gamma}(v)
&=\frac2{\pi hm}\int_{\mathbb R}
\frac{M_\Gamma(\omega)^2-M_\gamma(\omega)^2}{\omega^2}
|V_v(\omega)|^2\,d\omega\ \ge0.
\end{align*}
The apparent zero-frequency singularity has a continuous limit.
No zero-mean assumption on $v$ is needed.
\end{theorem}
The multiplier difference is the explicit gamma gain; $|V_v|^2$ describes
the tested pattern's frequencies; $1/\omega^2$ comes from underlying
step features. The last term accounts for finite geometry. Its sign and
size depend on direction, so a positive Fourier contribution alone does
not prove monotonicity of every finite-kernel eigenvalue or ratio.

\subsection*{Follow one direction the target actually needs}
At gamma 8, select the eigenvector with largest target energy among
resolved positive modes satisfying $\rho_i\le2\times10^{-6}$.
This reproducible retrospective selection uses the initial target and
kernel, with no training outcomes. It selects \textbf{eigenvector 22},
carrying \textbf{4.14\% of target squared norm}, with relative eigenvalue
only $1.612\times10^{-6}$.

Freeze this vector $v$ and change gamma. Starting from its gamma-8
response, the explicit mechanism predicts
\[
q^{\rm Fourier}_\Gamma(v)=q_8(v)+\mathcal A_{\Gamma,8}(v).
\]
This uses the known reference direction and gamma multiplier, with no
new eigenspace calculation at other slopes. Compare it with the directly
measured $\|J_\Gamma^Tv\|^2$. Numerically, the Fourier term is evaluated
through its equivalent closed-form pairwise entries
$2[d\coth(8d)-d\coth(\Gamma d)]/(hm)$, proved below.
Response rises from $0.000389636$ at gamma 8
to $0.386738$ at gamma 64, about a thousandfold. At gammas 12, 16, and
64, finite corrections are respectively \textbf{0.0601\%, 0.0339\%, and
0.00439\%} of actual response. On this target-relevant direction, the
explicit Fourier contribution quantitatively explains the change.

The direction stays fixed. At larger gamma, its response is a weighted
average of new eigenvalues, not necessarily one eigenvalue. Small measured
corrections apply to this probe and geometry. The full measured spectrum
below shows how change is distributed among actual learning directions.

\begin{figure}[htbp]\centering
\includegraphics[width=\linewidth]{optimizer_access_mechanism.pdf}
\caption{\textbf{Gamma's action is explicit and measurable.}
Left: analytic squared multipliers at all four gammas; marked frequencies
are the target's three nominal frequencies, not automatically finite-kernel
eigenmodes. Right: actual response and reference response plus Fourier
gain for fixed gamma-8 eigenvector 22, carrying 4.14\% initial target
energy. Curves nearly coincide because finite correction is below
0.061\% of actual response at every larger gamma. These are directional
responses, not eigenvalues tracked across gamma.}
\end{figure}
\clearpage

\section{The measured spectrum explains the GD delay}
For zero initialization, $r_n=J_\gamma w_n-y$ obeys
\[
r_{n+1}=(I-\eta K_\gamma)r_n.
\]
Expanding $y$ in orthonormal eigenvectors gives each component the factor
$1-\eta\mu_i$. Orthogonality yields exact relative squared error:
\[
e(n)^2=\sum_i p_i(\gamma)(1-\eta\mu_i(\gamma))^{2n}.
\]
Nullspace energy is included with factor one. The largest eigenvalue
limits the shared stable step: $0<\eta<2/\mu_1$. Our runs use
$\eta\mu_1\simeq1/2$, giving, at the exact normalized step,
\[
e(n)^2=\sum_i p_i(\gamma)\left(1-\frac{\rho_i(\gamma)}2\right)^{2n}.
\]
Forecasts use the actual archived step. A ratio of one halves a component's
residual each update; a ratio of $10^{-6}$ takes about two million updates
for one e-fold reduction. Target weights determine which rates matter.
Large eigenvalues elsewhere cannot accelerate an independently evolving
slow GD component.

For the fixed direction above, measured response divided by
$\mu_1(\gamma)$ is $1.612\times10^{-6}$ at gamma 8 and
$1.560\times10^{-3}$ at gamma 64, about a 968-fold increase.
This is the direction's weighted average of relative eigenvalues.
The complete target-weighted spectrum accounts for how the individual
rates contribute to training.

We measure the finite tanh spectrum and target weights, then evaluate this
classical formula without fitting a rate to training trajectories. First
1\% crossings are \textbf{15,798,313; 186,057; 61,792; and 16,013 updates}
at gammas 8, 12, 16, and 64. All four agree with executed crossings.
The gamma mechanism explains the kernel-response change; the spectral
calculation accounts for complete finite geometry when computing time.

\section{Does Adam encounter the same difficult directions?}
Adam rescales readout-coordinate updates and uses gradient history:
\[
r_{n+1}=r_n-\alpha_nJ_\gamma D_n\widehat m_n,\qquad
D_n=\diag\!\left((\sqrt{\widehat v_n}+\epsilon_{\rm Adam})^{-1}\right).
\]
$\widehat m_n$ and $\widehat v_n$ are bias-corrected first and second
moments. Even without momentum, the effective kernel changes with $D_n$.
GD eigenvectors are diagnostic coordinates for Adam; GD decay is not
an Adam time law.

We can still ask: \textbf{does error remain in directions the raw kernel
learns slowly?} Project saved Adam residuals onto the finite-kernel modes:
\[
\sum_{0<\rho_i(\gamma)\le\rho}
\frac{|u_i(\gamma)^Tr_n|^2}{\|y\|^2}.
\]
\begin{landscape}
\begin{figure}[p]\centering
\includegraphics[width=\linewidth]{optimizer_access_three_panel.pdf}
\caption{\textbf{From gamma to learning difficulty.}
A: leading 64 actual finite-kernel eigenvalue ratios, ordered separately
at each gamma; matching rank does not imply matching eigenvectors.
B: initial target energy and residual energy after 200,000 common-recipe
Adam updates, accumulated over positive modes below each ratio cutoff.
Energy $10^{-4}$ is the 1\% error budget. C: executed GD crossings and
measured-spectrum forecasts, with observed first Adam loss-EMA crossings
using one 100-update half-life. The arrow means no crossing by 200,000;
GD has longer budgets. Target and feature geometry match across panels.}
\end{figure}
\end{landscape}

All energies use original target norm. At gamma 8, modes with
$\rho_i\le2\times10^{-6}$ initially hold 4.33\% of target energy.
After 200,000 Adam updates, remaining energy is $1.47094\times10^{-4}$,
alone above the $10^{-4}$ budget for 1\% error. Almost all final residual
energy lies there. Adam removes about 99.66\% of initial band energy,
but enough remains to prevent acquisition.

At gammas 12, 16, and 64, final band energies are $2.2660\times10^{-6}$,
$1.0437\times10^{-7}$, and $3.8307\times10^{-8}$, below tolerance.
Subspaces share a spectral definition across gamma, while eigenvectors
may differ. Adam can temporarily increase projected residual components;
these measurements describe remaining error, not untouched initial content.

\subsection*{Make the Adam acquisition statistic visible}
The common setting uses initial rate $10^{-3}$, epsilon $10^{-12}$, and
moments $(0.9,0.999)$. Rate stays constant through 20,000 updates,
decays by cosine to $10^{-6}$ at 50,000, then remains fixed through
200,000. This shared schedule permits an early-learning comparison while
explaining later settling behavior.

Smooth squared relative error with shared half-life $s=100$ updates:
\[
M_0=e(0)^2=1,\qquad M_n=\beta M_{n-1}+(1-\beta)e(n)^2,
\qquad\beta=2^{-1/s}.
\]
The first update with $M_n\le10^{-4}$ is first EMA acquisition.
Plots show $\sqrt{M_n}$ in relative-error units. Common-recipe crossings
are \textbf{7,688; 1,900; and 1,486} for gammas 12, 16, and 64;
gamma 8 never crosses within the horizon. Every crossing is computed
from every saved update.

\begin{figure}[htbp]\centering
\includegraphics[width=\linewidth]{optimizer_access_schedule.pdf}
\caption{\textbf{The crossing comes from actual training.}
Bands show minima and maxima of raw errors in update bins; lines show
the square root of loss EMA. Dashed vertical lines and threshold dots
identify first crossings. The gray window is shared rate decay from
20,000 to 50,000. Displayed curves are sampled; crossing calculations
use every update.}
\end{figure}

This measures first acquisition of smoothed loss. Later EMA upcrossings
occur 73, 97, and 98 times at the three successful slopes. Raw sustained
crossings occur around 41,000 under the shared schedule. The EMA window
is a retrospective visualization choice with a memory-induced resolution
limit: $M_n\ge\beta^n$ precludes crossing before
$\lceil s\log_2(10^4)\rceil$. At half-life 100 this floor is
\textbf{1,329 updates}, close to the gamma-64 crossing. EMA compresses
differences between sufficiently fast runs.

For this attainable target and fixed geometry, small gamma leaves needed
directions weakly coupled to features. GD exhibits the delay implied by
its spectrum; Adam retains sufficient error in those slow directions to
miss the same accuracy threshold at gamma 8 within its budget.
\FloatBarrier

\section{Proof of the gamma-response theorem}
Let $I=[c_1-h/2,c_W+h/2]$ be the union of center cells, and put
\[
F_{ab,\gamma}(c)=\tanh(\gamma(x_a-c))\tanh(\gamma(x_b-c)).
\]
First, the whole-line deficit is
\[
\int_{\mathbb R}(1-F_{ab,\gamma}(c))\,dc
=2d_{ab}\coth(\gamma d_{ab}).
\]
For distinct samples, use
$1-\tanh u\tanh v=\coth(u-v)(\tanh u-\tanh v)$.
The translated-profile difference integrates to twice its translation.
For coincident samples, $\int\sech^2(\gamma c)\,dc=2/\gamma$ gives
the continuous limit.

Since $|I|=Wh$, subtract the deficit and restore its part outside $I$
to obtain $K_\gamma=B_\gamma+R_\gamma$, with
\begin{align*}
R_\gamma[a,b]
={}&\frac1m\left[\sum_{j=1}^{W}F_{ab,\gamma}(c_j)
-\frac1h\int_I F_{ab,\gamma}(c)\,dc\right]\\
&+\frac1{hm}\int_{\mathbb R\setminus I}(1-F_{ab,\gamma}(c))\,dc.
\end{align*}
The bracket is the quadrature discrepancy; the last integral accounts
for the finite interval. Bias contributes $1/m$ to $B_\gamma$ and
cancels between gammas.

Next, $s_\gamma*\sign=\tanh(\gamma\,\cdot)$ follows by differentiation:
both derivatives equal $2s_\gamma$ and both functions vanish at zero.
With $\widehat g(\omega)=\int g(x)e^{-i\omega x}\,dx$, substitute
$u=e^{2\gamma x}$ and $a=\omega/(2\gamma)$:
\begin{align*}
\widehat s_\gamma(\omega)
&=\int_0^\infty\frac{u^{-ia}}{(1+u)^2}\,du
=\Gamma(1-ia)\Gamma(1+ia)\\
&=\frac{\pi a}{\sinh(\pi a)}=M_\gamma(\omega).
\end{align*}
These equalities use the beta integral and gamma-function reflection
identity. Here $\Gamma(\cdot)$ is the gamma function, distinct from
the larger slope $\Gamma$.

Define the integrable difference profile
\[
D(d)=2[d\coth(\gamma d)-d\coth(\Gamma d)].
\]
Differentiating the deficit integral twice and integrating once by parts
gives
\[
\frac{d^2}{dd^2}[2d\coth(\gamma d)]=4(s_\gamma*s_\gamma)(d).
\]
To see every step, write the deficit as
$\int[1-t_\gamma(u)t_\gamma(u-d)]\,du$, with
$t_\gamma(u)=\tanh(\gamma u)$. Its second derivative is
\[
-\int t_\gamma(u)t_\gamma''(u-d)\,du
=\int t_\gamma'(u)t_\gamma'(u-d)\,du.
\]
The boundary term vanishes, and $t_\gamma'=2s_\gamma$, proving the identity.
Consequently
\[
\widehat D(\omega)=\frac{4(M_\Gamma(\omega)^2-M_\gamma(\omega)^2)}{\omega^2},
\qquad\widehat D(0)=\frac{\pi^2}{3}(\gamma^{-2}-\Gamma^{-2}).
\]
This follows from $\widehat{D''}=-\omega^2\widehat D$; expand
$z/\sinh z$ to get the continuous zero-frequency value. Since
$z/\sinh z$ decreases for positive $z$, $M_\Gamma\ge M_\gamma$,
so the transform is nonnegative.

Finally $B_\Gamma[a,b]-B_\gamma[a,b]=D(x_a-x_b)/(hm)$.
Fourier inversion contributes $1/(2\pi)$, and summation against
$v_av_b$ yields $|V_v(\omega)|^2$. This proves the Fourier term with
prefactor $2/(\pi hm)$. Add the exact finite correction to prove the
theorem.\hfill$\square$

\section{Empirical checks and reproduction}
The five targets are the sine mixture, $\exp(\sin(3\pi x))$,
$1/(1+25x^2)$, $\sqrt5x^2$, and $\sqrt2\sin(2\pi x)$.
Supplementary comparisons retain shared and archived validation-selected
Adam settings. Selection used median validation error at five checkpoints
from 40,000 to 50,000, over five initial rates and two epsilons;
optimizer state continued intact. These plots measure training-grid
optimization, without a new held-out generalization claim.

\begin{figure}[htbp]\centering
\includegraphics[width=\linewidth]{optimizer_access_targets.pdf}
\caption{\textbf{Target and optimizer setting matter.} All archived targets
use the same 100-update loss-EMA half-life. Selected settings optimize
late validation error, not early acquisition. Arrows indicate no crossing
within 200,000. Results do not yield universal monotone ordering by gamma.}
\end{figure}

The primary common-recipe ordering appears at EMA half-lives 30, 100,
and 300. At 1,000, crossings cluster near 28,000--30,000 and reorder.
With selected settings, half-life-100 crossings at gammas 12, 16, and
64 are 40,208; 41,716; and 10,076. These expose dependence on the
measurement window and optimizer recipe.

\begin{figure}[htbp]\centering
\includegraphics[width=\linewidth]{optimizer_access_ema_sensitivity.pdf}
\caption{\textbf{The smoothing choice is visible.} Each horizontal position
uses one half-life for all gammas. The dotted line marks 100 updates;
the gray dashed line is earliest possible crossing due to initial-loss
memory. Triangles mark censoring at 200,000. All points summarize the
same saved trajectories.}
\end{figure}
\FloatBarrier

All calculations use archived FP64 features and trajectories. Modal
analysis retains singular values above $10^{-14}$ of the largest;
omitted numerical directions are reported separately, not identified as
exact nullspace. Independent SVD and direct residual projections
reproduce slow-band energies within $5.1\times10^{-14}$. Every-update
traces determine EMA crossings; saved readouts determine modal residual
measurements. The fixed-direction calculation compares the explicit
increment to direct feature action and records the signed finite
correction. No new training was performed.

The \href{../results/checkpoint_D_optimizers/expD36_frozen_gamma_probe/full_sweep/refinements/gamma_optimizer_access/}{numerical evidence}
contains the selected direction, plotted quantities, source hashes, and
verification records. Reproduce from the repository root with one BLAS thread:
\par\vspace{6pt}
\begin{minipage}{\linewidth}\begin{quote}\footnotesize\begin{verbatim}
export OPENBLAS_NUM_THREADS=1 VECLIB_MAXIMUM_THREADS=1
python -m experiments.expD36_frozen_gamma_probe.gamma_mechanism_analysis
python -m experiments.expD36_frozen_gamma_probe.adam_spectral_analysis
python -m experiments.expD36_frozen_gamma_probe.adam_ema_analysis \
  --archive /path/to/adam_raw_scalar_archive.npz
python -m experiments.expD36_frozen_gamma_probe.optimizer_access_figure
latexmk -pdf -outdir=/tmp/gamma-optimizer-latex \
  docs/gamma_optimizer_access_note.tex
\end{verbatim}\end{quote}\end{minipage}
\par\vspace{6pt}
Extract the scalar archive from persistent training data using
\texttt{adam\_trace\_audit.py}. Existing compact EMA artifacts suffice
to regenerate figures without retrieving complete traces.
\end{document}
